\begin{table}[htbp]
  \centering
  \caption[¿Se estorban la estrategia y el tamaño?]{\textcolor{borrador2}{Acierto medio en el
      primer resultado de cada estrategia en cada tamaño máximo, promediando las tres
      configuraciones que la rejilla dedica a cada combinación. Cincuenta preguntas del
      conjunto de evaluación.}}
  \label{tab:interaccion_fragmentacion}
  \small
  \begin{tabular}{@{}rrrrrr@{}}
    \toprule
    \textbf{Máx.} & \textbf{Estructural} & \textbf{Tamaño fijo} & \textbf{Semántica} & \textbf{Rango} & \textbf{Truncan} \\
    \midrule
    600           & 0,930                & 0,930                & 0,930              & 0,000          & 0/9              \\
    \textbf{900}  & \textbf{0,930}       & 0,918                & 0,923              & \textbf{0,012} & \textbf{0/9}     \\
    1.200         & 0,863                & 0,883                & 0,883              & 0,020          & 0/9              \\
    1.500         & 0,820                & 0,857                & 0,837              & 0,037          & 7/9              \\
    1.800         & 0,793                & 0,843                & 0,823              & 0,050          & 8/9              \\
    \bottomrule
  \end{tabular}
  \begin{minipage}{0.95\textwidth}
    \vspace{4pt}
    \tiny
    \textbf{Notas:} cada celda es la media de \textbf{RU@1} ---proporción de preguntas cuya
    asignatura correcta sale en el primer resultado--- de las tres configuraciones que esa
    estrategia tiene con ese tamaño máximo, de modo que la comparación entre columnas se hace
    siempre con el mismo número de intentos por estrategia. \textbf{Rango}: diferencia entre
    la mejor y la peor estrategia de esa fila; una diferencia de 0,020 equivale a una pregunta
    de las cincuenta. \textbf{Truncan}: cuántas de las nueve configuraciones de esa fila
    producen algún fragmento que supera la ventana del modelo de incrustaciones y que este
    recorta en silencio. En negrita, el tamaño adoptado.
    \textcolor{borrador2}{La columna \textbf{Rango} crece con el tamaño: los dos factores no
      son independientes en toda la rejilla, pero sí en los tamaños que se pueden elegir.}
  \end{minipage}
\end{table}
